\title{Ticket Tailor}

\maketitle
\section*{Project Description}
TicketTailor is a social media application, event organiser, and notifier that lets you stay up to date on upcoming social events at the university. 
The application also serves as a social media platform, allowing you to follow clubs and societies to stay up to date on their events.

\section*{Project Deliverables}
An MVP will require the following functionalities:
\begin{enumerate}
\item A secure authentication system allowing users to maintain individual profiles containing descriptive metadata and club affiliation details.
\item An interactive map interface that visualises event locations based on geographic filtering (distance from a set coordinate) and user-defined interest categories.
\item A RSVP mechanism for members, with the system dynamically updating and persisting the aggregate attendee count.
\item Tools for committee members to initialise events with configurable parameters, including pricing (free vs. ticketed) and tiered visibility (private vs. public).
\item Functionality to manage the transition of events from "Club-Only" to "Public" status, triggered by defined time-offsets and supported by automated alerts to notify the broader user base.
\item Save the event to your calendar
\item A management interface for organisers to modify event details, coupled with a push-based notification engine to inform all registered attendees of changes in real-time.
\end{enumerate}

\section*{Quality Attributes}
Below are the main quality attributes the application will focus on. 
In addition to these, implement strong privacy and security measures to protect user data.
\begin{enumerate}
\item Scalability: The system is evaluated by its ability to maintain performance metrics such as map marker rendering and notification delivery speeds during periods of high demand, such as when a popular private event is made public. Evaluation will show how infrastructure resources scale in response to request spikes.
\item Interoperability: The system’s interoperability is evaluated by its ability to exchange and synchronise event data with external calendar platforms (e.g., Google Calendar, Outlook, or Apple Calendar). Evaluation will involve testing the integrity of the data handshake and ensuring that exported metadata, including event titles, geographic coordinates, and temporal data, is accurately parsed by third-party systems.
\item Reliability: The system is considered reliable if the event and attendee databases are replicated across multiple instances, ensuring that RSVP counts and event status changes are synced correctly. Evaluation involves demonstrating that users view the most current data even during high-frequency update periods (e.g., right before a major event).
\end{enumerate}

\end{document}